\section{几何 Satake 的范畴迹}
\label{S: cattrSat}

本节转到仿射情形。基本想法依旧，但技术细节更繁复。

设 $F$ 是剩余域 $k=\bF_q$ 有限的局部域，$q$-Frobenius 记作 $\sigma$。
回忆 \S \ref{S: affGrass} 中的记号，特别是 \eqref{ramified Witt vector} 与 \eqref{discs}。我们假设 $G$ 是 $\mO$ 上的连通约化群。为简便，记 $K=L^+G$。

由上一节的讨论可见，要研究 $p$ 进群的表示论，重要的是研究 $LG$ 上 $\Ad LG$-等变层或 $\Ad_\sigma LG$-等变层的范畴。然而仿射情形的困难之一在于：商 $Y=\frac{LG}{\Ad LG}$ 或 $Y=\frac{LG}{\Ad_\sigma LG}$ 远非代数叠，因此如何理解其上的层范畴并不清楚。我们的办法（应当等价于 \cite{Ga} 中概述的办法）是用 \S \ref{SS: cat of corr} 中引入的范畴 $\on{D}^\Hk(X)$ 代替 $\on{D}(Y)$，其中 $X=\frac{LG}{\Ad K}$ 或 $X=\frac{LG}{\Ad_\sigma  K}$，而 $\Hk=\Hk(X):=X\times_YX$\footnote{更确切地说，需要按注 \ref{R: DC enh} (2) 所述把 $\on{D}(Y)$ 定义为 $\Shv^{\Hk}(X)$ 的同伦范畴。但本文不需要这个更精细的版本。}。我们将会看到，虽然 $X$ 与 $\Hk$ 都不是代数叠，它们可以被有限维代数叠逼近，因此 $\on{D}^{\Hk}(X)$ 有意义。用 $\on{D}^\Hk(X)$ 代替 $\on{D}(Y)$ 的另一个好处是：前者经由上同调对应实现，与整体模空间（Shimura 簇或整体 Shtuka 的模空间）相关。我们还要指出，在 Fargues–Scholze 处理 $p$ 进群局部 Langlands 对应的办法中，他们用 Fargues–Fontaine 曲线上 $G$-丛的模空间上的层范畴替代所寻求的范畴 $\on{D}(\frac{LG}{\Ad_\sigma LG})$。它应当经由近循环函子与我们的范畴 $\on{D}^\Hk(X)$ 相关。不过，作者并不清楚如何在 Fargues–Scholze 的办法中触及 $\on{D}(\frac{LG}{\Ad LG})$。

一旦定义了合适的范畴，我们就可以取几何 Satake 对应的范畴迹，这将是下一节算术应用的关键要素。本节只讨论 $Y=\frac{LG}{\Ad_\sigma LG}$ 的情形；不过几乎所有构造也适用于 $Y=\frac{LG}{\Ad LG}$ 的情形，见注 \ref{R:affine springer}。

\subsection{局部 Shtuka 的模空间}

由于
$X=\frac{LG}{\Ad_\sigma K}$、$Y=\frac{LG}{\Ad_\sigma LG}$ 与 $X\times_YX$ 都不是代数叠，要把先前的讨论推广到仿射情形，第一步是使 \S \ref{SS: cat of corr} 中定义的范畴在此情形下有确切含义。结果表明，此时的 $\frac{LG}{\Ad_\sigma K}$ 可以解释为局部 Shtuka 的模空间。

设 $R$ 是完美 $k$-代数。若 $\mE$ 是 $D_R$ 上的 $G$-主丛，则它沿 $q$-Frobenius $\sigma:D_R\to D_R$ 的拉回记作 ${^\sigma}\mE$。
$\Spec R$ 上的局部 $r$ 重 $G$-Shtuka 是指 $D_{R}$ 上一列 $G$-主丛 $\mE_1,\ldots, \mE_r$，连同一串修正
\[\mE_r\dashrightarrow \mE_{r-1}\dashrightarrow\cdots\dashrightarrow \mE_0:={^\sigma}\mE_r.\]
若
\[\inv(\mE_i\dashrightarrow \mE_{i-1})\leq \mu_i,\]
则称其奇点被 $\mmu=(\mu_1,\ldots,\mu_r)$ 所界。我们把局部 $r$ 重 $G$-Shtuka 的模空间 $\on{Sht}^{r,\loc}$ 定义为如下预叠：它对每个 $R$ 赋予 $\Spec R$ 上局部 $r$ 重 $G$-Shtuka 的群胚。它是闭子预叠 $\on{Sht}^{\loc}_{\mmu}$ 之并（$\mmu$ 取遍），后者由奇点被 $\mmu=(\mu_1,\ldots,\mu_r)$ 所界的局部 Shtuka 组成。特别地，当 $r=1$ 时，$\on{Sht}^{r,\loc}$ 记作 $\Sht^\loc$，称为局部 $G$-Shtuka 的模空间，它是闭子预叠 $\Sht^\loc_\mu$ 之并。

由于 $W(R)$ 上每个 $G$-主丛在 $R$ 上 \'etale 局部可平凡化（例如见 \cite[引理 1.3]{Zh1}），有自然同构
\begin{equation}
\label{R:pre of Sht}
\Sht^\loc\cong \frac{LG}{\Ad_\sigma K},\quad \mbox{它限制为同构 } \Sht^\loc_\mu\cong \frac{\Gr_\mu^{(\infty)}}{\Ad_\sigma K},\ \ \ \mu\in \xcoch(T)/W.
\end{equation}
具体地说，记 $\Sht^{\loc,\Box}$ 为 $\Sht^\loc$ 上的 $K$-主丛，它对局部 $G$-Shtuka $(\beta:\mE\dashrightarrow {^\sigma}\mE)$ 连同平凡化 $\epsilon: \mE\simeq \mE^0$ 分类。则复合 $\mE^0\xrightarrow{\epsilon^{-1}}\mE\stackrel{\beta}{\dashrightarrow}{^\sigma}\mE\xrightarrow{\sigma(\epsilon)}\mE^0$ 定义元素 $g\in LG$，并诱导同构 $\Sht^{\loc,\Box}\cong LG$；在该同构下，$K$ 在 $\Sht^{\loc,\Box}$ 上的作用等同于 $\Ad_\sigma$ 在 $LG$ 上的作用。此外，若 $\inv(\beta)\leq \mu$，则 $g\in \Gr_{\mu}^{(\infty)}$。取商即得上述同构。

\begin{ex}\label{no singularity}
若 $\mmu=\mu=0$，则 $\Sht^{\loc}_{0}=\mathbf B G(\mO)$ 是投射有限群 $G(\mO)$ 的分类叠。
\end{ex}

\begin{ex}\label{loc Sht and pdiv}
若 $\mO=\bZ_p$，则 $R$ 上奇点被余权 $\omega_i$（如例 \ref{minu Sch} 中所定义）所界的 $\GL_n$-Shtuka 范畴，等价于 $R$ 上维数 $i$、高度 $n$ 的 $p$-可除群的范畴。该函子把 $p$-可除群 $\bX$ 送到它的 Dieudonné 模，而 Gabber 的一个定理断言这是等价。注意来自 $p$-可除群的 $\GL_n$-Shtuka 非常特殊：它们的奇点必须被极小余权所界。
\end{ex}

\begin{rmk}
我们刚才定义的是奇点在闭点 $s\in D$ 处的局部 Shtuka。也可以定义奇点在一般点 $\eta\in D$ 处的局部 Shtuka，甚至让奇点沿 $D$ 移动。
在混合特征情形，奇点沿 $D$ 的局部 Shtuka 与 Breuil–Kisin 模密切相关（\cite{SW}）。
\end{rmk}

\begin{rmk}
在等特征情形，可以定义整体 $G$-Shtuka，其中 $D$ 被有限域上的整体代数曲线替代。这正是 Drinfeld 最初的发明，它极大地推广了椭圆模的概念。不过目前尚不清楚数域情形下是否存在整体 $G$-Shtuka 的类似物。一些猜测可参见 \cite{S}。
\end{rmk}

同构 \eqref{R:pre of Sht} 显然推广为同构
$$\frac{LG\times^KLG\times^K\cdots\times^KLG}{\Ad_\sigma K}\cong \on{Sht}^{r,\loc},$$
在此同构下，与 \eqref{E: psigma} 类似的映射也有模空间的解释，通常称为\emph{部分 Frobenius 映射}。
固定如上的一列支配余权 $\mmu=(\mu_1,\ldots,\mu_r)$，并令 $\mu_0=\sigma(\mu_r)$。则部分 Frobenius 映射定义为
\begin{equation*}
\label{hpf}
\xymatrix@R=0pt{
F_\mmu\colon \on{Sht}_{(\mu_1,\ldots,\mu_r)}^{r,\loc} \ar[r] & \on{Sht}^{r,\loc}_{(\sigma(\mu_{r}),\mu_1,\ldots,\mu_{r-1})} \\
(\mE_r\stackrel{\beta_r}{\dashrightarrow} \mE_{r-1}\stackrel{\beta_{r-1}}{\dashrightarrow}\cdots\stackrel{\beta_2}{\dashrightarrow} \mE_1\stackrel{\beta_1}{\dashrightarrow} {^\sigma}\mE_r) \ar@{|->}[r] & (\mE_{r-1}\stackrel{\beta_2}{\dashrightarrow}\cdots\stackrel{\beta_r}{\dashrightarrow} \mE_1\stackrel{\beta_1}{\dashrightarrow} {^\sigma}\mE_r\stackrel{\sigma(\beta_r)}{\dashrightarrow} {^\sigma}\mE_{r-1}),
}
\end{equation*}
它是预叠的同构（即完美 $k$-代数上函子的同构）。

我们定义如下预叠的对应
\begin{equation}
\label{E;Hk Shtuka}
\xymatrix{&\Hk^{s,t}(\Sht^\loc)\ar_{\overleftarrow{h}^\loc}[dl]\ar^{\overrightarrow{h}^\loc}[dr]&\\
\Sht^{s,\loc}&&\Sht^{t,\loc},
}
\end{equation}
其中 $\Hk^{s,t}(\Sht^\loc)$ 是如下预叠：它对每个完美 $k$-代数 $R$，分类 $D_R$ 上 $G$-主丛修正的如下交换图
\begin{equation*}\label{qi for diff type}
\xymatrix@C=35pt{
\mE'_s \ar@{-->}[r]^{\beta'_s} \ar@{-->}[d]_\beta & \cdots  \ar@{-->}[r]^{\beta'_2} & \calE'_1 \ar@{-->}[r]^{\beta'_1} & {}^\sigma\mE'_s \ar@{-->}[d]^{\sigma(\beta)}\\
\mE_t \ar@{-->}[r]^{\beta_t} & \cdots  \ar@{-->}[r]^{\beta_2} & \calE_1 \ar@{-->}[r]^{\beta_1} & {}^\sigma\mE_t,
}
\end{equation*}
使得顶行（相应地底行）给出 $\Sht^{s,\loc}$（相应地 $\Sht^{t,\loc}$）的 $R$-点。
我们把这样的图称为从顶行到底行的\emph{Hecke 对应}。当 $s=t=1$ 时，把 $\Hk^{s,t}(\Sht^\loc)$ 简记作 $\Hk(\Sht^\loc)$。
于是与 \eqref{R:pre of Sht} 类似，有典范同构
$$\Hk(\Sht^\loc)\cong \frac{LG}{\Ad_\sigma K}\times_{\frac{LG}{\Ad_\sigma LG}} \frac{LG}{\Ad_\sigma K},$$
其中 $\overleftarrow{h}^\loc$ 与 $\overrightarrow{h}^\loc$ 是自然投影。
对两列支配余权 $\la_\bullet$ 与 $\mmu$，令
\[\Sht^{\loc}_{\la_\bullet\mid\mmu}:=(\overleftarrow{h}^\loc)^{-1}(\Sht^\loc_{\la_\bullet})\cap (\overrightarrow{h}^\loc)^{-1}(\Sht^\loc_{\mmu}),\]
有时把 $\overleftarrow{h}^\loc$ 与 $\overrightarrow{h}^\loc$ 在 $\Sht^{\loc}_{\la_\bullet\mid\mmu}$ 上的限制记作 $\overleftarrow{h}^\loc_{\la_\bullet}$ 与 $\overrightarrow{h}^\loc_\mmu$。
若进一步要求 $\beta$ 的相对位置 $\leq \nu$，就得到 $\Sht_{\la_\bullet|\mmu}^{\loc}$ 的闭子预叠 $\Sht_{\la_\bullet|\mmu}^{\nu,\loc}$。

\begin{rmk}
\label{R:fiber of hloc}
注意态射 $\overleftarrow{h}^\loc: \Hk(\Sht^\loc)\to\Sht^\loc$ 的纤维正是仿射 Grassmann 簇。对每个 $\mu$，态射
$\overleftarrow{h}^\loc: (\overrightarrow{h}^\loc)^{-1}(\Sht^{\loc}_\mu)\to \Sht^\loc$ 的纤维是仿射 Deligne–Lusztig 簇（定义可参见 \cite[\S 3.1]{Zh1}）。
\end{rmk}
